% Compile on Windows: powershell -ExecutionPolicy Bypass -File build_theory.ps1
\documentclass[11pt]{article}

\usepackage[margin=1in]{geometry}
\usepackage{amsmath,amssymb,amsthm}
\usepackage{booktabs,tabularx}
\usepackage{bm}
\usepackage[hidelinks]{hyperref}

\newtheorem{assumption}{Assumption}
\newtheorem{proposition}{Proposition}

\newcommand{\q}{\bm{q}}
\newcommand{\qd}{\dot{\bm{q}}}
\newcommand{\qdd}{\ddot{\bm{q}}}
\newcommand{\R}{\mathbb{R}}
\newcommand{\M}{\bm{M}}
\newcommand{\Lo}{\bm{L}}
\newcommand{\E}{\mathbb{E}}
\newcommand{\epsacc}{\varepsilon_{\mathrm{acc}}}
\newcommand{\dacc}{\delta_{\mathrm{acc}}}
\newcommand{\Ccov}{C_{\mathrm{cov}}}

\title{Physics-Structured Dynamics Models for Model-Based RL\\
\large Lagrangian formulation, network parameterization, and the controller}
\author{lagrangian-mbrl-franka}
\date{\today}

\begin{document}
\maketitle

\begin{abstract}
We learn the forward dynamics of the 7-DoF Franka Emika Panda,
$(\q,\qd,\bm{\tau})\mapsto\qdd$, to use as the predictive model inside a
model-based reinforcement-learning (MBRL) loop. This note states the
Euler--Lagrange equations the dynamics obey, shows \emph{exactly where the
neural network enters} a Deep Lagrangian Network (DeLaN), lists the architectures
used for the structured model and the unstructured baselines, fixes the
statistical error convention used by the project, and specifies the controller
that plans through the learned model.
\end{abstract}

\section{Problem}
A rigid-body manipulator with generalized coordinates $\q\in\R^{n}$
($n=7$ for the Franka) and applied joint torques $\bm{\tau}\in\R^{n}$ evolves
according to the manipulator equation
\begin{equation}
\M(\q)\,\qdd + \bm{c}(\q,\qd) + \bm{g}(\q) = \bm{\tau},
\label{eq:manip}
\end{equation}
Here $\M(\q)\succ0$ is the symmetric positive-definite mass matrix.
The vector $\bm{c}$ collects Coriolis and centrifugal terms, while $\bm{g}$ is
the gravitational force. The learning target is the \emph{forward dynamics}
\begin{equation}
\qdd = \M(\q)^{-1}\!\left(\bm{\tau}-\bm{c}(\q,\qd)-\bm{g}(\q)\right).
\label{eq:fwd}
\end{equation}

\section{Lagrangian mechanics}
The dynamics \eqref{eq:manip} are the Euler--Lagrange equations of the Lagrangian
$\mathcal{L}=T-V$, with kinetic energy $T=\tfrac12\qd^{\top}\M(\q)\qd$ and
potential $V(\q)$:
\begin{equation}
\frac{d}{dt}\!\left(\frac{\partial \mathcal{L}}{\partial \qd}\right)
-\frac{\partial \mathcal{L}}{\partial \q} \;=\; \bm{\tau}.
\label{eq:el}
\end{equation}
Carrying out the differentiation recovers \eqref{eq:manip} with
$\bm{g}(\q)=\partial V/\partial\q$ and the Coriolis term expressed through
derivatives of $\M$. Encoding this structure is the whole point: the true
dynamics \emph{lie in} this class, so a model constrained to it has a smaller
effective hypothesis space and needs fewer samples to reach a given accuracy
(the sample-complexity argument of \texttt{PROJECT\_PLAN.md}~\S2).

\section{Where the network goes: Deep Lagrangian Networks (DeLaN)}
\label{sec:delan}
DeLaN \cite{lutter2019delan} does \emph{not} regress $\qdd$ directly. Instead,
two neural networks parameterize the \emph{energy}, and the dynamics follow by
automatic differentiation of \eqref{eq:el}. Concretely, with parameters
$\bm{\theta}$:
\begin{itemize}
  \item a network outputs the entries of a lower-triangular
  $\Lo(\q;\bm{\theta})$, giving a guaranteed-SPD mass matrix
  \begin{equation}
  \M(\q;\bm{\theta}) = \Lo(\q;\bm{\theta})\,\Lo(\q;\bm{\theta})^{\top}
  + \varepsilon\bm{I}, \qquad \varepsilon>0,
  \label{eq:cholesky}
  \end{equation}
  (the diagonal of $\Lo$ is passed through a shifted softplus so it stays
  strictly positive, which conditions the solves);
  \item a scalar network outputs the potential $V(\q;\bm{\theta})$.
\end{itemize}
The kinetic-plus-potential Lagrangian
$\mathcal{L}=\tfrac12\qd^{\top}\M(\q;\bm\theta)\qd-V(\q;\bm\theta)$ is then
differentiated through \eqref{eq:el}. Using the identity implemented in the code,
the combined Coriolis-plus-gravity force is
\begin{equation}
\bm{c}(\q,\qd)+\bm{g}(\q) \;=\;
\left(\frac{\partial^{2}\mathcal{L}}{\partial\qd\,\partial\q}\right)\!\qd
\;-\;\frac{\partial \mathcal{L}}{\partial \q},
\label{eq:cg}
\end{equation}
and the predicted acceleration follows from \eqref{eq:fwd} by solving the SPD
system $\M\,\qdd = \bm\tau-\bm c-\bm g$ (a Cholesky solve). Because
\eqref{eq:cg} takes \emph{second} derivatives of $\mathcal{L}$, the network
activations must be twice continuously differentiable ($C^{2}$): we use
\textbf{softplus} (or \textbf{tanh}), never ReLU.

\paragraph{Training.}
We minimize a one-step loss on $(\q,\qd,\bm\tau,\qdd)$ transitions combining the
forward-dynamics residual (predicted vs.\ measured $\qdd$) and the
inverse-dynamics residual (predicted vs.\ measured $\bm\tau$), with input
normalization fit on the training set.

\section{Statistical target and finite-horizon guarantee}
\label{sec:bound}
The statistical target is the held-out \emph{one-step acceleration MSE}
\begin{equation}
  \epsacc(\hat f;\rho)
  := \E_{(\bm{x},\bm{\tau})\sim\rho}
  \left[\left\|\hat f(\bm{x},\bm{\tau})
  - f^\star(\bm{x},\bm{\tau})\right\|_2^2\right],
  \qquad \bm{x}=(\q,\qd).
  \label{eq:acc-mse}
\end{equation}
We write $\dacc:=\sqrt{\epsacc}$ for its root mean-squared error.  The
$H$-step rollout MSE remains an important diagnostic, but it is not the
statistical primitive in the bound: it already contains closed-loop
distribution shift and compounding error, which the simulation argument below
handles explicitly.

The code uses semi-implicit Euler with step $\Delta t$:
\[
  \qd^+ = \qd+\Delta t\,f(\bm{x},\bm{\tau}),\qquad
  \q^+ = \q+\Delta t\,\qd^+.
\]
Consequently an acceleration error $\bm e=\hat f-f^\star$ induces the exact
one-step state-transition error
\begin{equation}
  \left\|\hat F(\bm{x},\bm{\tau})-F^\star(\bm{x},\bm{\tau})\right\|_2
  = c_{\Delta t}\|\bm e\|_2,\qquad
  c_{\Delta t}:=\Delta t\sqrt{1+\Delta t^2}.
  \label{eq:integration-error}
\end{equation}

\begin{assumption}[Coverage and Lipschitz values]
\label{ass:coverage}
For every policy used by the optimality comparison and every time
$t\in\{0,\ldots,H-1\}$, the relevant state--action occupancy $d_t^\pi$ is
absolutely continuous with respect to the aggregated training distribution
$\rho$, with density ratio at most $\Ccov$.  The continuation value at time
$t+1$ is $L_{t+1}$-Lipschitz in the Euclidean state norm.
\end{assumption}

\begin{proposition}[MSE-to-control conversion]
\label{prop:simulation}
Let $\hat\pi$ be $\varepsilon_{\mathrm{opt}}$-optimal for the learned transition
$\hat F$, and let $\pi^\star$ be optimal for $F^\star$. Under
Assumption~\ref{ass:coverage},
\begin{equation}
  J_{F^\star}(\pi^\star)-J_{F^\star}(\hat\pi)
  \le
  2c_{\Delta t}\!\left(\sum_{t=0}^{H-1}L_{t+1}\right)
  \sqrt{\Ccov\,\epsacc(\hat f;\rho)}
  +\varepsilon_{\mathrm{opt}}.
  \label{eq:main-bound}
\end{equation}
\end{proposition}

\paragraph{Proof sketch.}
Apply the finite-horizon value-difference identity to a fixed policy and bound
each transition perturbation by the Lipschitz constant of the continuation
value.  Change measure from $d_t^\pi$ to $\rho$ using the density-ratio bound,
then apply Cauchy--Schwarz and \eqref{eq:integration-error}.  Applying the
fixed-policy bound once to $\pi^\star$ and once to $\hat\pi$, and inserting the
learned-model planning error, gives \eqref{eq:main-bound}. If
$L_{t+1}\le L_V(H-t-1)$, then the leading constant is at most
$c_{\Delta t}L_VH(H-1)$, displaying the usual $O(H^2)$ amplification.

\paragraph{Resolution of the square root.}
Equation~\eqref{eq:main-bound} is proportional to
$\sqrt{\epsacc}$ because \eqref{eq:acc-mse} is a \emph{squared} error.  The same
bound is linear in the RMSE $\dacc$.  We therefore report MSE in experiments,
take its square root only when inserting it into a control-performance bound,
and never use the same symbol for both quantities.

\subsection{On-policy data are not treated as i.i.d.}
Data are collected in rounds.  Conditional on the history before round $m$,
the policy $\pi_m$ is frozen and its Markov chain is assumed stationary after
burn-in and geometrically $\beta$-mixing:
\[
  \beta_m(k)\le B_m\exp(-k/\tau_m).
\]
For a block length $b$, alternating blocks yield
$N_{\mathrm{eff}}=\lfloor N/(2b)\rfloor$ approximately independent blocks; $b$
is chosen so the coupling remainder
$2N_{\mathrm{eff}}\beta_m(b)$ is below the allocated failure probability.
Supervised generalization is applied conditionally within each round and then
union-bounded over rounds.  The training distribution is the data-aggregation
mixture $\rho=\sum_m\alpha_m d^{\pi_m}$.  Assumption~\ref{ass:coverage} is still
required for the final policy; without it, the theorem makes no claim about
unvisited regions.

\paragraph{Claim scope.}
The paper's current claim is a \emph{conditional improvement of an upper
bound}: under realizability, coverage, mixing, and comparable optimization,
the structured class can have a smaller estimation-complexity term.  This is
not yet a two-sided separation.  A separation would additionally require a
matching lower bound for the unstructured class, which is outside the present
result and will not be implied by empirical curves alone.

\section{Model architectures (what we use in each case)}
\label{sec:arch}
\begin{table}[h]
\centering
\begin{tabularx}{\textwidth}{@{}p{0.17\textwidth}p{0.22\textwidth}X@{}}
\toprule
Role & Model & Architecture \\
\midrule
Structured (\emph{proposed}) & DeLaN \cite{lutter2019delan} &
  energy heads $\Lo(\q),V(\q)$, MLP $[128,128]$, \textbf{softplus} ($C^2$),
  SPD via \eqref{eq:cholesky} \\
Unstructured baseline & MLP ensemble (PETS) \cite{chua2018pets} &
  $K{=}5$ probabilistic MLPs $[256,256,256]$, SiLU, Gaussian $\qdd$ head \\
Minimal reference & MLP & single deterministic MLP $[256,256,256]$, SiLU \\
\bottomrule
\end{tabularx}
\caption{Models compared. The structured model constrains the hypothesis class
to physically consistent dynamics; the ensemble is the strong unstructured
baseline and supplies epistemic uncertainty for planning. Widths/activations are
the defaults; the benchmark tunes them per model.}
\end{table}

\paragraph{Related structured models (future ablations).}
Lagrangian Neural Networks \cite{cranmer2020lnn} learn an \emph{arbitrary}
$\mathcal{L}$ (dropping the $T=\tfrac12\qd^\top\M\qd$ assumption); Hamiltonian
Neural Networks \cite{greydanus2019hnn} parameterize the Hamiltonian instead;
and Lagrangian Graph Neural Networks \cite{bhattoo2022lgnn} exploit the
manipulator's kinematic topology and scale to many links. They slot into the
same model registry as drop-in alternatives to DeLaN.

\section{The controller: planning through the learned model}
\label{sec:control}
Given the learned forward model $\hat f_{\bm\theta}:(\q,\qd,\bm\tau)\mapsto\qdd$
and a one-step integrator $\hat F$ producing $(\q_{t+1},\qd_{t+1})$, control is
\emph{model-predictive}. At each state $\bm{x}_t=(\q_t,\qd_t)$ we optimize a
torque sequence over a horizon $H$,
\begin{equation}
\min_{\bm\tau_{t:t+H-1}} \;
\sum_{k=0}^{H-1}\ell(\bm{x}_{t+k},\bm\tau_{t+k}) + \ell_f(\bm{x}_{t+H}),
\qquad
\bm{x}_{t+k+1}=\hat F(\bm{x}_{t+k},\bm\tau_{t+k}),
\label{eq:mpc}
\end{equation}
apply the first torque, and re-plan (receding horizon). The default solver is
\textbf{MPPI} \cite{williams2017mppi}: sample $N$ noisy control sequences, roll
each through $\hat F$, and update the control by the cost-weighted average
\begin{equation}
\bm\tau \;\leftarrow\; \sum_{i=1}^{N} w_i\,\bm\tau^{(i)},
\qquad
w_i \propto \exp\!\Big(-\tfrac1\lambda\,\big(J^{(i)}-\min_j J^{(j)}\big)\Big),
\label{eq:mppi}
\end{equation}
with $J^{(i)}$ the rolled-out cost \eqref{eq:mpc} and temperature $\lambda$.
MPPI needs \emph{no policy network}, so its performance is governed almost
entirely by the accuracy of $\hat f_{\bm\theta}$---this is the cleanest probe of
whether the Lagrangian prior helps. The reach cost is the end-effector position
(and orientation) error to the commanded pose plus a torque-effort penalty.

\paragraph{Alternatives.}
A Dyna/MBPO variant \cite{janner2019mbpo} would instead generate short model
rollouts and run an off-policy update (e.g.\ SAC) on mixed real+model data;
a model-free PPO policy serves as the lower-bound baseline. These are deferred
until MPPI results justify a learned policy.

\section{Evaluation protocol}
Models are compared on held-out one-step acceleration MSE (and, for the
structured model, energy drift and $\lambda_{\min}(\M)$), under matched data:
the same seeds and the same per-seed transitions feed every model. We tune
hyper-parameters per model, retrain the winner over $\ge 5$ seeds, and report the
mean with a percentile-bootstrap 95\% confidence interval. The data-efficiency
story is the curve of model error (and, later, task return) versus the number of
real transitions $N$.

\begin{thebibliography}{9}
\bibitem{lutter2019delan} M.~Lutter, C.~Ritter, J.~Peters.
  \emph{Deep Lagrangian Networks: Using Physics as Model Prior for Deep
  Learning.} ICLR, 2019.
\bibitem{cranmer2020lnn} M.~Cranmer et al.
  \emph{Lagrangian Neural Networks.} ICLR Deep Differential Equations
  Workshop, 2020.
\bibitem{greydanus2019hnn} S.~Greydanus, M.~Dzamba, J.~Yosinski.
  \emph{Hamiltonian Neural Networks.} NeurIPS, 2019.
\bibitem{bhattoo2022lgnn} R.~Bhattoo, S.~Ranu, N.~M.~A.~Krishnan.
  \emph{Learning Articulated Rigid Body Dynamics with Lagrangian Graph Neural
  Network.} NeurIPS, 2022.
\bibitem{chua2018pets} K.~Chua, R.~Calandra, R.~McAllister, S.~Levine.
  \emph{Deep Reinforcement Learning in a Handful of Trials using Probabilistic
  Dynamics Models (PETS).} NeurIPS, 2018.
\bibitem{janner2019mbpo} M.~Janner, J.~Fu, M.~Zhang, S.~Levine.
  \emph{When to Trust Your Model: Model-Based Policy Optimization (MBPO).}
  NeurIPS, 2019.
\bibitem{williams2017mppi} G.~Williams, A.~Aldrich, E.~A.~Theodorou.
  \emph{Model Predictive Path Integral Control.} Journal of Guidance, Control,
  and Dynamics, 2017.
\bibitem{agarwal2022rltheory} A.~Agarwal, N.~Jiang, S.~M.~Kakade, W.~Sun.
  \emph{Reinforcement Learning: Theory and Algorithms.} Online monograph, 2022.
\bibitem{mohri2010mixing} M.~Mohri, A.~Rostamizadeh.
  \emph{Stability Bounds for Stationary $\phi$-mixing and $\beta$-mixing
  Processes.} Journal of Machine Learning Research, 2010.
\end{thebibliography}

\end{document}
